\subsection{Evaluation questions and scope}\label{subsec:eval-questions-and-scope}

Each of the four research questions maps to a different measurement procedure. Table~\ref{tab:research-questions} summarizes this mapping, and the mapping is detailed in the following sections.

\begin{table}[h]
\centering
    \begin{tabular}{@{}lll@{}}
        \toprule
        \textbf{RQ} & \textbf{Object measured} & \textbf{Defining section} \\
        \midrule
        RQ1 & Rubric-scored document-extraction faithfulness
        & Section~3.6.3 \\
        RQ2 & Knowledge-extraction strict F1 and provenance carry-rate
        & Section~3.6.4 \\
        RQ3 & Cascade cost--quality trade-off
        & Section~3.6.5 \\
        RQ4 & NLI behaviour in grounding and relation classification
        & Section~3.6.6 \\
        \bottomrule
    \end{tabular}
    \caption{Research questions and their measurement procedures.}
    \label{tab:research-questions}
\end{table}

This section does not report empirical findings; the purpose of this section is to define the evaluation design, in order to make the results in Chapter~\ref{chapter:Results} interpretable: what is measured, the metrics that were used for storing, how the datasets and the cases were sampled, and the offline-replay protocol used for cascade calibration and evaluation. 